\section*{Problem 1}

反向传播算法的核心思想是先计算高层误差，再将误差逐层传回低层。设一个两层线性网络为
\[
\bm{h} = W_1 \bm{x},\quad \hat{\bm{y}} = W_2 \bm{h},
\]
其中
\[
W_1 = \begin{bmatrix}
    1 & 2 \\
    0 & -1
\end{bmatrix},\quad
W_2 = \begin{bmatrix}
    1 & 0 \\
    2 & -1
\end{bmatrix},\quad
\bm{x} = \begin{pmatrix} 1 \\ -1 \end{pmatrix},\quad
\bm{y} = \begin{pmatrix} 0 \\ 1 \end{pmatrix}.
\]
损失函数为
\[
L = \frac{1}{2} \|\hat{\bm{y}} - \bm{y}\|_2^2.
\]
完成以下任务：
\begin{enumerate}
    \item 计算 $\bm{h}$、$\hat{\bm{y}}$、$\bm{e} = \hat{\bm{y}} - \bm{y}$。
    \item 计算 $\nabla_{W_2} L$。
    \item 计算 $\nabla_{W_1} L$。
\end{enumerate}

\begin{proof}
    \begin{enumerate}
        \item 前向计算得到
        \[
        \begin{aligned}
            \bm{h}
            &=W_1\bm{x}
             =\boxed{\begin{pmatrix}-1\\1\end{pmatrix}},\\
            \hat{\bm{y}}
            &=W_2\bm{h}
             =\boxed{\begin{pmatrix}-1\\-3\end{pmatrix}},\\
            \bm{e}
            &=\hat{\bm{y}}-\bm{y}
             =\boxed{\begin{pmatrix}-1\\-4\end{pmatrix}}.
        \end{aligned}
        \]
        \item 因为 $\nabla_{\hat{\bm{y}}}L=\bm{e}$，所以
        \[
            \nabla_{W_2}L
            =\bm{e}\bm{h}^\top
            =\boxed{\begin{bmatrix}1&-1\\4&-4\end{bmatrix}}.
        \]
        \item 传回隐藏层的梯度为
        \[
            \nabla_{\bm{h}}L
            =W_2^\top\bm{e}
            =\begin{pmatrix}-9\\4\end{pmatrix}.
        \]
        因而
        \[
            \nabla_{W_1}L
            =(W_2^\top\bm{e})\bm{x}^\top
            =\boxed{\begin{bmatrix}-9&9\\4&-4\end{bmatrix}}.
        \]
    \end{enumerate}
\end{proof}

\section*{Problem 2}

自动微分可以看作是系统地应用链式法则。设计算图为
\[
L = \bm{c}^\top \bm{v},\quad
\bm{v} = \tanh(\bm{u}),\quad
\bm{u} = A \bm{x} + \bm{b},
\]
其中
\[
A = \begin{bmatrix}
    1 & -1 \\
    0 & 2
\end{bmatrix},\quad
\bm{x} = \begin{pmatrix} 1 \\ 0 \end{pmatrix},\quad
\bm{b} = \begin{pmatrix} -1 \\ 0 \end{pmatrix},\quad
\bm{c} = \begin{pmatrix} 2 \\ -1 \end{pmatrix}.
\]
这里 $\tanh$ 逐坐标作用于向量。完成以下任务：
\begin{enumerate}
    \item 进行前向计算，求 $\bm{u}$、$\bm{v}$、$L$。
    \item 使用反向模式自动微分思想，求 $\nabla_{\bm{x}} L$、$\nabla_{A} L$、$\nabla_{\bm{b}} L$。
\end{enumerate}

\begin{proof}
    \begin{enumerate}
        \item 前向计算得到
        \[
            \bm{u}=A\bm{x}+\bm{b}
            =\begin{pmatrix}0\\0\end{pmatrix},\qquad
            \bm{v}=\tanh(\bm{u})
            =\begin{pmatrix}0\\0\end{pmatrix},\qquad
            L=\bm{c}^\top\bm{v}=\boxed{0}.
        \]
        \item 从输出向输入反向计算。因为 $\tanh'(0)=1$，
        \[
            \nabla_{\bm{v}}L=\bm{c},\qquad
            \nabla_{\bm{u}}L
            =\bm{c}\odot\bigl(\bm{1}-\bm{v}\odot\bm{v}\bigr)
            =\begin{pmatrix}2\\-1\end{pmatrix}.
        \]
        因而
        \[
        \begin{aligned}
            \nabla_{\bm{x}}L
            &=A^\top\nabla_{\bm{u}}L
             =\boxed{\begin{pmatrix}2\\-4\end{pmatrix}},\\
            \nabla_AL
            &=(\nabla_{\bm{u}}L)\bm{x}^\top
             =\boxed{\begin{bmatrix}2&0\\-1&0\end{bmatrix}},\\
            \nabla_{\bm{b}}L
            &=\nabla_{\bm{u}}L
             =\boxed{\begin{pmatrix}2\\-1\end{pmatrix}}.
        \end{aligned}
        \]
    \end{enumerate}
\end{proof}

\section*{Problem 3}

一维卷积或互相关可以看作是一个具有特殊稀疏结构的矩阵乘法。设一维输入信号为
\[
\bm{x} = (x_1, x_2, x_3, x_4, x_5)^\top,
\]
滤波器为
\[
\bm{a} = (a_1, a_2, a_3)^\top.
\]
采用步长为 $1$、无零填充的窄互相关运算：
\[
y_i = a_1 x_i + a_2 x_{i+1} + a_3 x_{i+2},\quad i = 1, 2, 3.
\]
完成以下任务：
\begin{enumerate}
    \item 将该运算写成 $\bm{y} = K \bm{x}$ 的矩阵形式。
    \item 令 $\bm{a} = (1, 0, -1)^\top$，$\bm{x} = (2, 1, 3, 0, -2)^\top$，计算 $\bm{y}$。
    \item 若目标输出为 $\bm{t} = (0, 1, 3)^\top$，损失函数为 $L = \frac{1}{2} \|\bm{y} - \bm{t}\|_2^2$，求 $\nabla_{\bm{x}} L$ 和 $\nabla_{\bm{a}} L$。
\end{enumerate}

\begin{proof}
    \begin{enumerate}
        \item 互相关矩阵为
        \[
            K=\begin{bmatrix}
                a_1&a_2&a_3&0&0\\
                0&a_1&a_2&a_3&0\\
                0&0&a_1&a_2&a_3
              \end{bmatrix},
            \qquad
            \boxed{\bm{y}=K\bm{x}}.
        \]
        \item 代入 $\bm{a}=(1,0,-1)^\top$，得到
        \[
            K=\begin{bmatrix}
                1&0&-1&0&0\\
                0&1&0&-1&0\\
                0&0&1&0&-1
              \end{bmatrix},
            \qquad
            \bm{y}=K\bm{x}
            =\boxed{\begin{pmatrix}-1\\1\\5\end{pmatrix}}.
        \]
        \item 记残差
        \[
            \bm{r}=\bm{y}-\bm{t}
            =\begin{pmatrix}-1\\0\\2\end{pmatrix}.
        \]
        由 $\bm{y}=K\bm{x}$，有
        \[
            \nabla_{\bm{x}}L
            =K^\top\bm{r}
            =\boxed{\begin{pmatrix}-1\\0\\3\\0\\-2\end{pmatrix}}.
        \]
        对滤波器系数，令
        \[
            X_{\mathrm{win}}
            =\begin{bmatrix}
                x_1&x_2&x_3\\
                x_2&x_3&x_4\\
                x_3&x_4&x_5
              \end{bmatrix}
            =\begin{bmatrix}
                2&1&3\\
                1&3&0\\
                3&0&-2
              \end{bmatrix},
        \]
        则 $\bm{y}=X_{\mathrm{win}}\bm{a}$，所以
        \[
            \nabla_{\bm{a}}L
            =X_{\mathrm{win}}^\top\bm{r}
            =\boxed{\begin{pmatrix}4\\-1\\-7\end{pmatrix}}.
        \]
    \end{enumerate}
\end{proof}
